\begin{abstract}
Vertical federated learning (VFL) lets organisations that hold different features of the same entities train a joint model without pooling their data; with images, for example, satellites operated by different organisations observe the same areas, or different clinics hold complementary images of the same skin lesion. The standard method, split learning, exchanges activations and gradients at every mini-batch, so every party must stay online for thousands of messages. FedOnce needs a single upload, but each party must first train its own encoder on its own, often small, data. We make two contributions. (C1)~We revisit FedOnce with public foundation models: each party encodes its images with a frozen pre-trained encoder, optionally compresses the features with its own PCA, and uploads them once; the label owner then trains the classifier alone. Under one accounting of bytes, messages, party compute and what the label owner sees, we compare this protocol with split learning, logit exchanges, differential privacy and encrypted training. It is the best default: on VertiBench Satellite, \frVbReachSplitNNFoundBytes{} per sample match the accuracy of split learning, which exchanges \frVbReachSplitNNFoundRatio{} times more, and adapting the encoder by split learning gains a few points only on dermatology images, at the price of keeping every party online. (C2)~We assemble benchmarks whose parties hold genuinely different views of the same entity rather than pieces of one image: VertiBench Satellite, reproduced and audited, and a more realistic variant built from the same WorldStrat source, whose two parties hold different sensors and whose folds are grouped by site; and MILK10k, which pairs a clinical and a dermoscopic photograph of each lesion.
\keywords{Vertical federated learning \and Foundation models \and One-shot learning \and Benchmarks \and Communication efficiency \and Privacy}
\end{abstract}
